\section{Topology: exact reduction and geometric obligations}

\subsection{What the exact arithmetic already proves}

Let
\[
p=12\ell_0-4\ell_1-3\ell_2.
\]
The retained exact harness verifies the determinant and Smith-normal-form calculations for the printed relation matrix. In particular, the manuscript's chosen parameters $(\ell_0,\ell_1,\ell_2)=(0,1,-1)$ give $p=-1$.

\begin{proposition}[terminal group calculation]\label[proposition]{prop:terminal-group}
\Verified\ For the printed based presentation and its stated meridian relations, the resulting group is cyclic of order $|p|$. Hence the printed presentation is trivial for $(0,1,-1)$.
\end{proposition}

\begin{proof}
The exact row and column operations recorded in the audit reduce the displayed
$2\times2$ relation matrix to Smith normal form
$\operatorname{diag}(1,|p|)$. A group presented by that matrix is therefore
$\mathbb Z/|p|\mathbb Z$, which is trivial when $p=-1$.
\end{proof}

\begin{remark}
Proposition~\ref{prop:terminal-group} is a theorem about the printed
presentation. It becomes a theorem about $\pi_1(X)$ only after the based
inclusions induced by all three analytic attaching maps have been derived in
the same generators and orientations.
\end{remark}

\subsection{The map that must be identified}

Write the relevant decomposition schematically as $X=U\cup N$, where $N$ is the union of the filling pieces and $U\cap N$ is the disjoint union of their collars. The homological boundary map is
\[
\delta_q=(i_{U*},-i_{N*})\colon
H_q(U\cap N;\mathbb Z)\longrightarrow
H_q(U;\mathbb Z)\oplus H_q(N;\mathbb Z).
\]
The printed matrices denoted $\alpha_q$ have verified ranks, Smith forms, and saturation properties. The central question is not another Smith-form computation; it is the equality
\[
\boxed{\alpha_q=[\delta_q]_{\mathcal B_q}}
\]
for explicitly constructed primitive oriented integral bases $\mathcal B_q$.

The equality is decomposed below into the six nonredundant proof obligations
\hyperref[obl:C1]{C1}, \hyperref[obl:C2]{C2}, \hyperref[obl:C3]{C3},
\hyperref[obl:V1]{V1}, \hyperref[obl:MV1]{MV1}, and
\hyperref[obl:L1]{L1}.  The first three construct the finite equivariant
model, compute its integral collapse map, and identify that map with
nearby-cycle specialization.  Obligation V1 fixes the based van Kampen
inclusions; MV1 identifies the actual Mayer--Vietoris boundary cylinders and
multiple-fibre indices; L1 supplies an independently oriented
Leray--Mayer--Vietoris comparison.  This crosswalk is exhaustive for the
conditional endpoint stated here; none of the six is replaced by matching a
rank or by choosing the desired matrix as a normalization.

\begin{test}[fast falsification targets]
Any of the following disproves the claimed topology calculation: a nonprimitive geometric image where $\alpha_q$ assumes a primitive one; an index or torsion class omitted by specialization; a meridian sign changing $p=-1$ to a nonunit; or a Leray transgression differing from $\pm p\omega$ by a nonunit factor.
\end{test}

\subsection{Conditional recognition}

\begin{theorem}[precise conditional endpoint]
\Conditional\ Assume obligations \hyperref[obl:C1]{C1}--\hyperref[obl:C3]{C3},
\hyperref[obl:V1]{V1}, \hyperref[obl:MV1]{MV1}, and
\hyperref[obl:L1]{L1} are discharged and give exactly the printed integral
maps. Then the verified terminal calculations yield
\[
\pi_1(X)=0,\qquad H_k(X;\mathbb Z)=
\begin{cases}
\mathbb Z,&k=0,6,\\
0,&1\le k\le5.
\end{cases}
\]
Consequently $X$ is a smooth homotopy $6$-sphere; the standard fact $\Theta_6=0$ identifies it diffeomorphically with $S^6$.
\end{theorem}

\begin{proof}
The hypotheses identify the actual van Kampen and Mayer--Vietoris/Leray maps with the already verified presentations and matrices. Their Smith forms give the displayed fundamental group and homology. A simply connected integral homology sphere is a homotopy sphere by Hurewicz and Whitehead; the final smooth recognition uses $\Theta_6=0$.
\end{proof}
